%%%%%%%%%%%%%%%%%%%%%%%%%%%%%%%%%
%
%       EXERCÍCIO
%
%%%%%%%%%%%%%%%%%%%%%%%%%%%%%%%%%

\definevalorquestao

\begin{exercicioBanco}[\valorquestao]
Uma das maneiras de saber se o seu peso está adequado à sua altura é calculando o Índice de Massa Corporal (IMC), que pode ser obtido através da fórmula dada por
%
\[
    h^{2}\cdot \text{IMC} - P = 0,
\]
%
sendo \(h\) a altura da pessoa em metros (m) e \(P\) o seu peso em quilogramas (kg). Se uma pessoa possui IMC igual a 30 kg/m\(^2\) e está com peso igual a 120 kg, indique a alternativa que apresenta o valor de sua altura em metros.

% Define as alternativas
\newcommand{\alternativas}{%
    \begin{center}
        \begin{tabularx}{\textwidth}{XXXXX}
            (a) \(1,7\). &
            (b) \(1,8\). &
            (c) \(1,9\). &
            (d) \(2,0\). &
            (e) \(2,1\).
        \end{tabularx}
    \end{center}
}

% Define a resposta correta
\newcommand{\resposta}{D}

% Lógica condicional para exibição
\ifdefstring{\atividade}{lista}{%
    \alternativas
    \vspace{0.5em}
    
    \noindent\textbf{Resposta:} letra \textbf{\resposta}.
}{%
    \ifdefstring{\modo}{objetivo}{%
        \alternativas
    }{%
        % modo = discursiva → não mostra alternativas
    }
}
\end{exercicioBanco}

